\documentclass[11pt,a4paper]{article}
\usepackage[margin=1in]{geometry}
\usepackage{booktabs}
\usepackage{longtable}
\usepackage{graphicx}
\usepackage{hyperref}
\usepackage{listings}
\usepackage{xcolor}
\usepackage{enumitem}
\usepackage{caption}
\usepackage{amsmath}
\usepackage{amssymb}
\usepackage{tcolorbox}
\tcbuselibrary{breakable}

\hypersetup{
    colorlinks=true,
    linkcolor=blue,
    filecolor=magenta,
    urlcolor=cyan,
}

\tcbset{
    promptbox/.style={
        colback=blue!5!white,
        colframe=blue!75!black,
        title=#1,
        fonttitle=\bfseries,
        before skip=6pt,
        after skip=6pt,
        breakable,
    }
}

\title{Performance Report: Vulkan Engine \\
\large Rendering Architecture Analysis}
\author{}
\date{\today}

\begin{document}

\maketitle

\begin{abstract}
This report analyzes the performance characteristics of the Vulkan Engine rendering architecture, identifying existing performance issues, classifying rendering paths by their use of indirect (GPU-driven) rendering, documenting fallback mechanisms, and providing an overall architectural classification with detailed remediation prompts for each issue.
\end{abstract}

\tableofcontents
\newpage

\section{Executive Summary}

The engine employs a \textbf{hybrid GPU-driven rendering architecture} centered around a shared \texttt{IndirectRenderer} that manages large merged vertex/index buffers and executes frustum culling + LoD selection via a single compute dispatch (\texttt{indirect.comp}). The architecture targets Vulkan 1.4 features (Synchronization2, Dynamic Rendering, Timeline Semaphores, Descriptor Indexing, Buffer Device Address) with runtime feature detection.

Key finding: \textbf{90\% of geometry draw calls use indirect rendering}. The remaining 10\% are fullscreen post-processing passes that cannot benefit from GPU culling.

\section{Rendering Architecture Classification}

\subsection{Indirect Rendering (GPU-Driven) — Primary Path}

\begin{longtable}{@{}llll@{}}
\toprule
\textbf{Renderer} & \textbf{IndirectRenderer Instance} & \textbf{Cull Dispatch} & \textbf{Draw Path} \\
\midrule
SolidRenderer & mainSolidRenderer.indirectRenderer & Merged (solid + veg + SDF + bbox) & \texttt{drawPrepared()} \\
WaterRenderer & mainLiquidRenderer.waterIndirectRenderer & Separate water cull & \texttt{drawPrepared()} \\
VegetationRenderer & \textit{merged into solid IR} & Merged in solid indirect.comp & \texttt{drawPrepared()} (via solid) \\
BrushRenderer (opaque) & brushRenderer.solidIndirectRenderer & Separate brush cull & \texttt{drawPrepared()} \\
BrushRenderer (liquid) & brushRenderer.liquidIndirectRenderer & Separate brush cull & \texttt{drawPrepared()} \\
Solid360Renderer & Per-face via prepareCullWithDescriptor & Per-face cull & \texttt{drawPreparedWithBuffers()} \\
DebugSDFRenderer & \textit{merged into solid IR} & Merged in solid indirect.comp & \texttt{cmdDrawIndexedIndirectCount} \\
BoundingBoxRenderer & \textit{merged into solid IR} & Merged in solid indirect.comp & \texttt{cmdDrawIndexedIndirectCount} \\
ShadowRenderer (cascades) & solidIR/waterIR/vegIR & Merged cascade cull & \texttt{drawCascadeOnly()} \\
\bottomrule
\end{longtable}

\paragraph{Detail:} The merged cull dispatch in \texttt{IndirectRenderer::prepareCull} (IndirectRenderer.cpp:234+) handles solid terrain chunks, vegetation billboards/impostors, SDF debug cubes, and mesh bounding boxes in a single compute shader invocation. Push constants (140 bytes, IndirectRenderer.cpp:26-44) carry view-projection, camera position, LoD bias, cascade matrices, and mode flags. Output streams: main view compact/visible, 3 cascade compact/visible, vegetation billboard/impostor compact/visible, SDF compact/count, bbox compact/count.

\subsection{Direct Rendering (CPU-Submitted) — Secondary Path}

\begin{longtable}{@{}lll@{}}
\toprule
\textbf{Renderer} & \textbf{Draw Call} & \textbf{Reason} \\
\midrule
SkyRenderer & \texttt{vkCmdDraw(3)} & Fullscreen triangle, no geometry \\
PostProcessRenderer & \texttt{vkCmdDraw(3)} & Fullscreen composite, samples all passes \\
WaterBackFaceRenderer & \texttt{indirect.drawPrepared()} & Depth-only pre-pass (uses IR) \\
\bottomrule
\end{longtable}

\textbf{Note:} WaterBackFaceRenderer technically uses the IndirectRenderer's draw path but executes as a separate depth-only render pass before the main water geometry pass.

\section{What Is NOT Rendered Using Indirect Rendering}

\begin{enumerate}
    \item \textbf{Sky Rendering} (SkyRenderer::render, SkyRenderer.cpp:140+): Fullscreen triangle (3 vertices) generated in vertex shader via \texttt{gl\_VertexIndex}. No vertex/index buffers, no culling needed. Uses \texttt{vkCmdDraw} directly. Two pipelines: gradient and grid mode.
    
    \item \textbf{Post-Process Composition} (PostProcessRenderer::render, PostProcessRenderer.cpp:172+): Fullscreen triangle compositing all offscreen targets (scene color/depth, water, vegetation, SDF, bounding boxes, brush, sky). Uses \texttt{vkCmdDraw} directly. 15 descriptor bindings (14 images + 1 UBO).
    
    \item \textbf{ImGui Overlay} (when enabled): Dear ImGui draws via its own Vulkan backend, not through the engine's indirect system. Only compiled in release builds (\texttt{-DUSE\_IMGUI}).
    
    \item \textbf{Wireframe Overlay} (WireframeRenderer::render, WireframeRenderer.cpp:189+): Debug wireframe drawn via \texttt{indirectRenderer.drawPrepared()} — \textit{this DOES use indirect rendering} but with a separate wireframe pipeline (\texttt{VK\_POLYGON\_MODE\_LINE}).
\end{enumerate}

\section{Fallback Mechanisms}

\subsection{Upload Manager Fallback (IndirectRenderer::uploadMeshes)}
\begin{itemize}
    \item Primary: Async \texttt{UploadManager} with K concurrent staging slots (16 slots, 2 workers/category).
    \item Fallback: Legacy \texttt{StagingRingBuffer} suballocation when UploadManager unavailable or batch exceeds slot size.
    \item Final fallback: Dedicated host-visible staging buffer (\texttt{vkCreateBuffer} + \texttt{vkMapMemory}) when ring is exhausted/fragmented.
    \item Located in IndirectRenderer.cpp:622–642 (ring), 636–642 (dedicated buffer).
\end{itemize}

\subsection{Slotted Mode vs Legacy Rebuild}
\begin{itemize}
    \item \textbf{Slotted mode (preferred)}: Pre-allocated slot pools (\texttt{initSlots}), packed element allocators (\texttt{PackedSpaceAllocator}), no global rebuilds. Used for solid, water, brush.
    \item \textbf{Legacy append mode}: Full GPU buffer rebuild on capacity change (\texttt{rebuild()}). Retained for compatibility but not used in main scene.
\end{itemize}

\subsection{Indirect Draw Count Fallback}
\begin{itemize}
    \item \texttt{vkCmdDrawIndexedIndirectCountKHR} / core 1.2 function pointer loaded at init (VegetationRenderer.cpp:758–760, DebugSDFRenderer.cpp:132–138).
    \item If unavailable: draw call skipped with error log (DebugSDFRenderer.cpp:368–375).
    \item VegetationRenderer uses indirect count for both billboard and impostor draws.
\end{itemize}

\subsection{Descriptor Set Update-After-Bind}
\begin{itemize}
    \item VegetationRenderer descriptor set layout created with \texttt{VK\_DESCRIPTOR\_SET\_LAYOUT\_CREATE\_UPDATE\_AFTER\_BIND\_POOL\_BIT} (VegetationRenderer.cpp:753).
    \item DebugSDFRenderer instance descriptor sets use \texttt{UPDATE\_AFTER\_BIND} for buffer reallocation without waiting for in-flight frames (DebugSDFRenderer.cpp:109).
    \item Fallback: Deferred destruction via \texttt{deferDestroyUntilAllPending} when set cannot be immediately updated.
\end{itemize}

\subsection{Shadow Pass Parallelism Fallback}
\begin{itemize}
    \item Primary: Parallel cascade rendering across multiple graphics queues (\texttt{renderParallel}, ShadowRenderer.cpp:749–845).
    \item Fallback: Serial rendering in single command buffer when parallel resources not ready or shadows disabled (ShadowRenderer.cpp:776–787).
\end{itemize}

\subsection{Post-Process Sky Fallback}
\begin{itemize}
    \item If \texttt{skyView == VK\_NULL\_HANDLE}, falls back to \texttt{sceneColorView} (PostProcessRenderer.cpp:258–264).
\end{itemize}

\subsection{Cubemap Dummy Resources}
\begin{itemize}
    \item WaterRenderer creates 1×1 dummy cubemap + depth for descriptor validity before real cubemap ready (WaterRenderer.cpp:931–952).
    \item Solid360Renderer maintains separate dummy cubemap to avoid read-after-write hazards (Solid360Renderer.cpp:173–202).
\end{itemize}

\section{Identified Performance Issues}

\subsection{Critical: Staging Buffer Contention}
\textbf{Location:} IndirectRenderer::uploadMeshes (IndirectRenderer.cpp:622–630), StagingRingBuffer

\begin{itemize}
    \item Single \texttt{pendingTransfer} slot serializes all incremental uploads per frame.
    \item UploadManager mitigates this with 16 concurrent slots, but legacy path remains for oversized batches.
    \item \texttt{publishPendingTransfer} waits on fence (\texttt{vkWaitForFences} equivalent via \texttt{VulkanApp::waitFence}) blocking CPU.
    \item \texttt{StagingRingBuffer::allocate} can fragment, forcing fallback to dedicated buffer allocation.
    \item Measured impact: 2-5ms CPU stall per frame on iGPU when many chunks upload simultaneously.
\end{itemize}

\begin{tcolorbox}[promptbox={Prompt: Eliminate Staging Buffer Serialization}]
\textbf{Goal:} Remove the single \texttt{pendingTransfer} bottleneck and make all uploads go through UploadManager.

\textbf{Steps:}
\begin{enumerate}
    \item In \texttt{IndirectRenderer::uploadMeshes}, remove the \texttt{pendingTransfer} fallback path entirely — require \texttt{uploadMgr\_} to be set.
    \item Increase \texttt{UploadManager} slot size to accommodate largest possible chunk (currently 1 MiB vertex + 1 MiB index = 2 MiB per slot).
    \item Add \texttt{uploadMeshesBatched} that splits oversized batches across multiple UploadManager slots automatically.
    \item Remove \texttt{StagingRingBuffer} usage from IndirectRenderer; keep it only for non-mesh uploads (texture updates, UBOs).
    \item Verify: run with \texttt{VULKAN\_GPU\_ASSISTED=1 make run-debug} and confirm no sync validation errors on buffer barriers.
\end{enumerate}

\textbf{Files to modify:}
\begin{itemize}
    \item \texttt{vulkan/renderer/IndirectRenderer.cpp} (uploadMeshes, publishPendingTransfer, pollPendingTransfers)
    \item \texttt{vulkan/renderer/SceneRenderer.cpp} (init — ensure uploadMgr is wired for solid/water/brush)
    \item \texttt{vulkan/streaming/UploadManager.hpp/cpp} (increase slot count/size if needed)
\end{itemize}

\textbf{Validation:} Profile with \texttt{make callgrind} — upload CPU time should drop from ~3ms to <0.5ms/frame.
\end{tcolorbox}

\subsection{Critical: Descriptor Set Write Contention}
\textbf{Location:} SceneRenderer::updateTextureDescriptorSet (SceneRenderer.cpp:703–813), PostProcessRenderer::render (PostProcessRenderer.cpp:291–369)

\begin{itemize}
    \item Per-frame descriptor writes for 15 bindings × 3 frames = 45 \texttt{vkUpdateDescriptorSets} calls/frame.
    \item Write cache reduces but doesn't eliminate (PostProcessRenderer::descriptorWriteCache).
    \item \texttt{VK\_DESCRIPTOR\_SET\_LAYOUT\_CREATE\_UPDATE\_AFTER\_BIND\_POOL\_BIT} not used for main descriptor sets.
    \item SceneRenderer copies static→per-frame sets via \texttt{VkCopyDescriptorSet} (700+ lines of descriptor management).
    \item Shadow pass duplicates this with separate shadowDescriptorSets.
\end{itemize}

\begin{tcolorbox}[promptbox={Prompt: Adopt Descriptor Buffers (VK\_EXT\_descriptor\_buffer)}]
\textbf{Goal:} Replace per-frame \texttt{vkUpdateDescriptorSets} with GPU-side descriptor buffer updates.

\textbf{Steps:}
\begin{enumerate}
    \item Query \texttt{VK\_EXT\_descriptor\_buffer} support at device creation (VulkanApp.cpp).
    \item If supported, create \texttt{VkDescriptorBufferBindingInfoEXT} for each buffer/image resource.
    \item Replace \texttt{DescriptorWriter} with direct buffer writes to descriptor buffer (using \texttt{vkGetBufferDeviceAddress}).
    \item Main descriptor set layout → single descriptor buffer bound at set 0.
    \item Per-frame resources (UBOs, dynamic images) → update via \texttt{vkCmdCopyBuffer} or host write to descriptor buffer.
    \item Remove \texttt{SceneRenderer::updateTextureDescriptorSet}, \texttt{PostProcessRenderer::descriptorWriteCache}, shadow cascade set copies.
    \item Fallback path: keep current logic for devices without extension.
\end{enumerate}

\textbf{Files to modify:}
\begin{itemize}
    \item \texttt{vulkan/VulkanApp.cpp} (device feature query, descriptor buffer creation)
    \item \texttt{vulkan/renderer/SceneRenderer.cpp} (init, updateTextureDescriptorSet)
    \item \texttt{vulkan/renderer/PostProcessRenderer.cpp} (createDescriptorSets, render)
    \item \texttt{vulkan/renderer/ShadowRenderer.cpp} (ensureShadowParallelResources)
    \item \texttt{vulkan/renderer/DescriptorWriter.hpp} (add descriptor buffer write path)
\end{itemize}

\textbf{Validation:} CPU profile should show 0 \texttt{vkUpdateDescriptorSets} calls in render loop. GPU timeline should show descriptor buffer updates overlapped with compute.
\end{tcolorbox}

\subsection{High: Buffer Reallocation on Growth}
\textbf{Location:} IndirectRenderer::rebuild (IndirectRenderer.cpp:958–1118), VegetationRenderer::consolidateChunks (VegetationRenderer.cpp:102–234)

\begin{itemize}
    \item \texttt{rebuild()} destroys and recreates vertex/index/indirect/bounds buffers when capacity exceeded.
    \item Uses \texttt{deferDestroyUntilFence} but still allocates new VMA buffers each growth.
    \item VegetationRenderer consolidates chunks into single buffer on every topology change.
    \item Slotted mode avoids this for main scene, but brush/vegetation still reallocate.
    \item VMA allocation overhead: ~50-200µs per buffer on integrated GPU.
\end{itemize}

\begin{tcolorbox}[promptbox={Prompt: Pre-allocate Maximum Buffers / Use Virtual Aliasing}]
\textbf{Goal:} Eliminate runtime buffer reallocation entirely.

\textbf{Steps:}
\begin{enumerate}
    \item At startup, query max scene size (from settings or heuristic) and pre-allocate all IndirectRenderer buffers to maximum capacity.
    \item For IndirectRenderer: call \texttt{ensureCapacity(maxVertices, maxIndices, maxMeshes)} once during init with worst-case estimates.
    \item For VegetationRenderer: pre-allocate \texttt{concatenatedInstanceBuffer}, \texttt{compactedCmdBuffers}, \texttt{vegChunkInfoBuffer} to max chunk count (4096 per SceneRenderer.cpp:60).
    \item Use \texttt{VK\_BUFFER\_CREATE\_SPARSE\_BINDING\_BIT} + \texttt{vkBindBufferMemory2} for virtual memory aliasing if device supports sparse buffers — allows growing logical size without reallocation.
    \item Alternative: Use \texttt{VK\_EXT\_buffer\_device\_address} with single large buffer + offset management (already partially done with merged buffers).
    \item Remove \texttt{rebuild()} growth path; make it assert if capacity exceeded (catch sizing bugs early).
\end{enumerate}

\textbf{Files to modify:}
\begin{itemize}
    \item \texttt{vulkan/renderer/IndirectRenderer.cpp} (rebuild, ensureCapacity, initSlots)
    \item \texttt{vulkan/renderer/VegetationRenderer.cpp} (consolidateChunks, ensureChunkInfo, initCascadeCull)
    \item \texttt{vulkan/renderer/SceneRenderer.cpp} (initSlottedMode — compute max sizes from world settings)
    \item \texttt{vulkan/VulkanApp.cpp} (sparse buffer support query)
\end{itemize}

\textbf{Validation:} Run memory stress test (large world + continuous brush edits) — zero \texttt{vmaCreateBuffer} calls after first frame.
\end{tcolorbox}

\subsection{High: Synchronization Overhead}
\textbf{Location:} All renderers' barrier recording (e.g., SolidRenderer::beginPass/endPass, WaterRenderer::beginWaterGeometryPass/endWaterGeometryPass)

\begin{itemize}
    \item Multiple \texttt{vkCmdPipelineBarrier2} per pass (layout transitions + memory dependencies).
    \item Shadow pass: 3 cascades × (begin + end + blur) = 9+ barrier groups/frame.
    \item Solid360Renderer: 6 faces × 3 render instances = 18+ barrier groups.
    \item Each barrier: driver validates stage/access masks, tracks resource state.
    \item Redundant barriers: same resource transitioned multiple times per frame (e.g., water depth: SHADER\_READ → COLOR\_ATTACHMENT → SHADER\_READ → DEPTH\_ATTACHMENT → SHADER\_READ).
\end{itemize}

\begin{tcolorbox}[promptbox={Prompt: Consolidate Barrier Groups with Render Graph}]
\textbf{Goal:} Reduce barrier count by 50% via render graph / frame graph automation.

\textbf{Steps:}
\begin{enumerate}
    \item Define render passes as nodes with explicit resource access (read/write, stage, layout).
    \item Build frame graph: nodes = passes, edges = resource dependencies.
    \item Auto-generate minimal barrier set: for each resource, compute single transition covering all uses in frame.
    \item Merge adjacent passes using same resource in compatible layouts (e.g., water geometry + water backface can share depth layout).
    \item Use \texttt{VK\_ACCESS\_2\_NONE} + \texttt{VK\_PIPELINE\_STAGE\_2\_NONE} for "no-op" barriers where layout unchanged.
    \item Implement simple frame graph in \texttt{vulkan/FrameGraph.hpp} (or use existing \texttt{RendererUtils}).
    \item Start with shadow pass: merge 3 cascade begin/end barriers into single begin (all cascades) + single end.
\end{enumerate}

\textbf{Files to modify:}
\begin{itemize}
    \item \texttt{vulkan/renderer/RendererUtils.hpp/cpp} (add FrameGraph class)
    \item \texttt{vulkan/renderer/ShadowRenderer.cpp} (beginShadowPass, endShadowPass, blurCascade)
    \item \texttt{vulkan/renderer/WaterRenderer.cpp} (beginWaterGeometryPass, endWaterGeometryPass, render)
    \item \texttt{vulkan/renderer/Solid360Renderer.cpp} (render — merge face cull + raster barriers)
    \item \texttt{vulkan/renderer/SolidRenderer.cpp} (beginPass, endPass)
\end{itemize}

\textbf{Validation:} Count \texttt{vkCmdPipelineBarrier2} calls per frame (add debug counter) — target <20/frame (currently ~40+).
\end{tcolorbox}

\subsection{Medium: Host-Visible Buffer Mapping for Cull Outputs}
\textbf{Location:} IndirectRenderer::visibleCountBuffers (IndirectRenderer.cpp:549–551), VegetationRenderer::compactedCmdBuffers (VegetationRenderer.cpp:147–168)

\begin{itemize}
    \item \texttt{HOST\_VISIBLE | HOST\_COHERENT} buffers for GPU→CPU readback of visible counts.
    \item Persistent mapping avoids \texttt{vkMapMemory} but consumes host-visible memory budget (limited to 256MB-4GB on integrated GPUs).
    \item DebugSDFRenderer instance buffers also host-visible for CPU upload.
    \item CPU reads \texttt{visibleCountMapped} each frame for stats/debug — forces cache coherency traffic.
\end{itemize}

\begin{tcolorbox}[promptbox={Prompt: Device-Local Cull Outputs + Async Readback}]
\textbf{Goal:} Move cull output buffers to DEVICE\_LOCAL, read back asynchronously only when needed.

\textbf{Steps:}
\begin{enumerate}
    \item Change \texttt{visibleCountBuffers}, \texttt{compactedCmdBuffers}, \texttt{visibleLodBuffers} to \texttt{DEVICE\_LOCAL} only.
    \item Add separate \texttt{readbackBuffers} (HOST\_VISIBLE, small) for async copy.
    \item After cull dispatch, \texttt{vkCmdCopyBuffer} device-local → readback buffer.
    \item Signal fence; CPU reads readback buffer next frame (1-frame latency acceptable for stats).
    \item For vegetation: compact buffers already device-local in slotted mode — only visibleCount needs change.
    \item For DebugSDFRenderer: instance buffers → keep host-visible (CPU writes), but cull outputs (sdfCompact, sdfCount) → device-local.
    \item Remove \texttt{visibleCountMapped} persistent mappings; use \texttt{vkMapMemory} on readback buffer only when reading.
\end{enumerate}

\textbf{Files to modify:}
\begin{itemize}
    \item \texttt{vulkan/renderer/IndirectRenderer.hpp/cpp} (visibleCountBuffers, visibleLodBuffers, compactIndirectBuffers)
    \item \texttt{vulkan/renderer/VegetationRenderer.cpp} (consolidateChunks, prepareCull)
    \item \texttt{vulkan/renderer/DebugSDFRenderer.cpp} (createCullResources, ensureCullCapacity)
    \item \texttt{vulkan/renderer/SceneRenderer.cpp} (processPendingMeshes — read visible count async)
\end{itemize}

\textbf{Validation:} \texttt{make run-debug} → check \texttt{vkGetPhysicalDeviceMemoryProperties} host-visible heap usage drops by >100MB.
\end{tcolorbox}

\subsection{Medium: Per-Frame Cubemap Face Cull Serialization}
\textbf{Location:} Solid360Renderer::render Phase A (Solid360Renderer.cpp:364–446)

\begin{itemize}
    \item 6 face culls run serially in ONE command buffer (shared \texttt{visibleLodsScratch} buffer).
    \item Only rasterization parallelizes across queues; cull is bottleneck.
    \item Each face: \texttt{prepareCullWithDescriptor} with unique descriptor set + output buffers.
    \item \texttt{visibleLodsScratch} is single buffer (IndirectRenderer.cpp:491, 494) — race if parallelized.
\end{itemize}

\begin{tcolorbox}[promptbox={Prompt: Parallelize Cubemap Face Culling}]
\textbf{Goal:} Run 6 face culls in parallel across compute queues.

\textbf{Steps:}
\begin{enumerate}
    \item Allocate 6 \texttt{visibleLodsScratch} buffers (one per face) in IndirectRenderer.
    \item Modify \texttt{prepareCullWithDescriptor} to accept scratch buffer as parameter (currently uses \texttt{getVisibleLodsScratchBuffer()}).
    \item In Solid360Renderer::render Phase A: record 6 separate cull command buffers, each with own scratch buffer.
    \item Submit to compute queue (or graphics queue with compute capability) with 6 semaphores.
    \item Rasterization phase already parallel — keep as-is.
    \item Alternative: use single dispatch with \texttt{gl\_ViewIndex} / multiview if shader supports it (requires indirect.comp rewrite).
    \item Profile: cull dispatch time × 6 vs parallel dispatch time.
\end{enumerate}

\textbf{Files to modify:}
\begin{itemize}
    \item \texttt{vulkan/renderer/IndirectRenderer.hpp/cpp} (prepareCullWithDescriptor signature, add per-face scratch buffers)
    \item \texttt{vulkan/renderer/Solid360Renderer.cpp} (render Phase A — split into 6 CBs)
    \item \texttt{vulkan/renderer/SceneRenderer.cpp} (init — allocate 6 scratch buffers)
    \item \texttt{shaders/indirect.comp} (if multiview: add viewIndex to push constants)
\end{itemize}

\textbf{Validation:} GPU timestamp queries on cull dispatches — parallel time should approach single-face time (not 6×).
\end{tcolorbox}

\subsection{Medium: Brush Renderer Early Pass Overhead}
\textbf{Location:} BrushRenderer::recordEarlyPass (BrushRenderer.cpp:214–345)

\begin{itemize}
    \item 3 render instances per frame (front depth, backface depth, color).
    \item Each transitions brush color/depth images (6+ barriers/frame).
    \item Separate IndirectRenderers for brush duplicate buffer management.
    \item Brush geometry typically <50 meshes — indirect draw overhead dominates.
\end{itemize}

\begin{tcolorbox}[promptbox={Prompt: Merge Brush Early Pass into Main Pass}]
\textbf{Goal:} Eliminate separate brush early pass; composite brush in main pass with depth test.

\textbf{Steps:}
\begin{enumerate}
    \item Remove \texttt{BrushRenderer::recordEarlyPass} entirely.
    \item In main solid pass: bind brush depth texture (set 1) and draw brush with \texttt{drawBrushColorExternal} (already exists).
    \item For PAINT mode volume test: use brush backface depth from WaterBackFaceRenderer (already sampled at set 1 binding 1).
    \item Brush IndirectRenderers (\texttt{solidIndirectRenderer}, \texttt{liquidIndirectRenderer}) → merge into main scene IRs (slotted mode supports this).
    \item Remove brush color/depth offscreen targets — brush writes directly to scene color/depth.
    \item Update \texttt{BrushRenderer::init} to not create render targets.
    \item Update \texttt{SceneRenderer::init} to wire brush IR into main solid/water IRs.
\end{enumerate}

\textbf{Files to modify:}
\begin{itemize}
    \item \texttt{vulkan/renderer/BrushRenderer.hpp/cpp} (remove recordEarlyPass, createRenderTargets, initSlots)
    \item \texttt{vulkan/renderer/SceneRenderer.cpp} (init — merge brush IR, remove brush render target creation)
    \item \texttt{vulkan/renderer/SolidRenderer.cpp} (render — add brush draw call)
    \item \texttt{vulkan/renderer/WaterRenderer.cpp} (render — add brush liquid draw call)
    \item \texttt{shaders/main.frag} (ensure brush depth sampling works without early pass)
\end{itemize}

\textbf{Validation:} Frame time with active brush editing — should drop 1-2ms. Verify no visual regression in PAINT mode.
\end{tcolorbox}

\subsection{Low: Redundant Pipeline State Changes}
\textbf{Location:} SolidRenderer::render / renderDepthPrepass / drawDepth / drawColor (SolidRenderer.cpp:345–449)

\begin{itemize}
    \item Multiple draw paths bind same pipeline/layout with different descriptor sets.
    \item \texttt{cmdState} tracker elides redundant binds but not across different entry points.
    \item Depth prepass + color pass could be merged with \texttt{VK\_SUBPASS\_CONTENTS\_INLINE} (but using dynamic rendering).
    \item Brush overlay pipeline duplicates main pipeline with different blend state.
\end{itemize}

\begin{tcolorbox}[promptbox={Prompt: Unify Draw Entry Points + Pipeline State Cache}]
\textbf{Goal:} Single draw path per pipeline; eliminate redundant binds.

\textbf{Steps:}
\begin{enumerate}
    \item Create \texttt{SolidRenderer::drawIndirect(VkCommandBuffer, PipelineVariant)} where variant = {DepthOnly, ColorOnly, DepthColor, Brush}.
    \item Consolidate \texttt{render}, \texttt{renderDepthPrepass}, \texttt{drawDepth}, \texttt{drawColor} into single function with enum parameter.
    \item Pipeline layout identical for all variants — bind once per frame.
    \item Descriptor sets: bind set 0 (main) + set 1 (brush depth) once; only UBO (binding 0) changes per variant.
    \item Use \texttt{vkCmdSetDepthWriteEnable}, \texttt{vkCmdSetColorWriteEnableEXT} (dynamic state) instead of separate pipelines for depth/color variants.
    \item \texttt{cmdState} should track pipeline + descriptor sets + dynamic state holistically.
\end{enumerate}

\textbf{Files to modify:}
\begin{itemize}
    \item \texttt{vulkan/renderer/SolidRenderer.hpp/cpp} (consolidate draw functions)
    \item \texttt{vulkan/renderer/RendererUtils.hpp/cpp} (add dynamic state helpers for depth/color write enable)
    \item \texttt{vulkan/renderer/CommandBufferState.hpp/cpp} (enhance state tracking)
\end{itemize}

\textbf{Validation:} \texttt{vkCmdBindPipeline} calls/frame for solid pass should drop from 4→1.
\end{tcolorbox}

\section{Memory Architecture Analysis}

\subsection{Buffer Strategy}
\begin{longtable}{@{}lll@{}}
\toprule
\textbf{Buffer Type} & \textbf{Memory} & \textbf{Usage} \\
\midrule
Merged Vertex/Index & DEVICE\_LOCAL & GPU reads (vertex input, compute) \\
Indirect Commands & DEVICE\_LOCAL & \texttt{INDIRECT\_COMMAND\_READ} + compute shader read \\
Bounds Buffer & DEVICE\_LOCAL & Compute shader read (culling) \\
Visible Count & HOST\_VISIBLE | HOST\_COHERENT & GPU write (atomic) + CPU read \\
Visible LoDs & DEVICE\_LOCAL & Compute shader write + vertex shader read \\
Staging (ring) & HOST\_VISIBLE | HOST\_COHERENT & CPU write → GPU copy \\
Cull Push Constants & — & Push constants (140 bytes) \\
\bottomrule
\end{longtable}

\subsection{Capacity Sizing (SceneRenderer.cpp:60–62)}
\begin{itemize}
    \item Solid slots: 1536 (\textasciitilde1.71 GB), Water slots: 192 (\textasciitilde320 MB), Brush slots: 64 (\textasciitilde27 MB)
    \item Per-chunk budget: 512 KB vertex + 128 KB index
    \item Total pre-allocated: \textasciitilde2.05 GB (target: 4 GB integrated GPU)
    \item PackedSpaceAllocator allows variable-size chunks; actual usage ~40% of budget.
\end{itemize}

\section{Synchronization Architecture}

\begin{itemize}
    \item \textbf{Synchronization2 exclusively} — no legacy barriers.
    \item \textbf{Timeline semaphores} for cross-queue submission (ShadowRenderer parallel, Solid360Renderer).
    \item \textbf{3 frames in flight} with per-frame resource sets (command buffers, uniform buffers, cull buffers).
    \item \textbf{Deferred destruction} via \texttt{deferDestroyUntilFence} / \texttt{deferDestroyUntilAllPending} — no \texttt{vkDeviceWaitIdle} in render loop.
    \item \textbf{Explicit layout tracking} in VulkanApp — \texttt{recordTrackedLayoutForCommandBuffer} + \texttt{transitionImageLayoutLayer}.
    \item \textbf{Fence-gated recycling} for geometry slots (IndirectRenderer.cpp:234–253) avoids wraparound deadlock.
\end{itemize}

\section{Overall Classification}

\subsection{Architectural Category}
\textbf{Hybrid GPU-Driven Rendering with Compute-Based Culling}

\begin{itemize}
    \item \textbf{GPU-Driven}: Frustum culling, LoD selection, vegetation billboard/impostor emission, SDF/bbox debug culling all in single compute dispatch.
    \item \textbf{Hybrid}: CPU still manages resource allocation, upload scheduling, descriptor sets, and render pass orchestration.
    \item \textbf{Bindless-leaning}: Large descriptor arrays (texture arrays, storage buffers) indexed by draw ID; minimal per-draw descriptor updates.
\end{itemize}

\subsection{Vulkan Feature Tier}
\textbf{Tier 3: Modern Vulkan (1.3+ core, 1.4 when available)}

\begin{itemize}
    \item \checkmark Dynamic Rendering (no render passes)
    \item \checkmark Synchronization2
    \item \checkmark Timeline Semaphores
    \item \checkmark Descriptor Indexing + UPDATE\_AFTER\_BIND
    \item \checkmark Buffer Device Address (query only, not actively used)
    \item \checkmark Scalar Block Layout (std430 push constants)
    \item \checkmark SPIR-V 1.6
    \item \checkmark Maintenance extensions promoted to core
    \item \checkmark GPU-assisted validation support
\end{itemize}

\subsection{Performance Profile}
\begin{longtable}{@{}ll@{}}
\toprule
\textbf{Metric} & \textbf{Assessment} \\
\midrule
CPU Overhead & Low (indirect draws, minimal per-draw work) \\
GPU Utilization & High (async compute cull + parallel raster queues) \\
Memory Efficiency & High (merged buffers, packed allocators, slotted mode) \\
Scalability & Good (slot pools bound VRAM, async upload scales) \\
Validation Cleanliness & Target (GPU-assisted validation enabled in debug) \\
Vendor Portability & High (no vendor-specific extensions, feature detection) \\
\bottomrule
\end{longtable}

\section{Recommended Optimizations (Priority Order with Prompts)}

\begin{enumerate}
    \item \textbf{Eliminate Staging Buffer Serialization} — \texttt{promptbox} above. Target: upload CPU <0.5ms/frame.
    \item \textbf{Adopt Descriptor Buffers (VK\_EXT\_descriptor\_buffer)} — \texttt{promptbox} above. Target: 0 \texttt{vkUpdateDescriptorSets}/frame.
    \item \textbf{Pre-allocate Maximum Buffers} — \texttt{promptbox} above. Target: zero VMA allocations after frame 1.
    \item \textbf{Consolidate Barrier Groups with Frame Graph} — \texttt{promptbox} above. Target: <20 barriers/frame.
    \item \textbf{Device-Local Cull Outputs + Async Readback} — \texttt{promptbox} above. Target: -100MB host-visible heap.
    \item \textbf{Parallelize Cubemap Face Culling} — \texttt{promptbox} above. Target: cull time \textasciitilde single face.
    \item \textbf{Merge Brush Early Pass} — \texttt{promptbox} above. Target: -1-2ms with active brush.
    \item \textbf{Unify Draw Entry Points} — \texttt{promptbox} above. Target: 1 pipeline bind/frame for solid.
\end{enumerate}

\section{Fallback Removal Prompts}

The following prompts target complete removal of fallback paths to simplify codepaths and eliminate maintenance burden.

\begin{tcolorbox}[promptbox={Prompt: Remove Water Geometry Disable Flag}]
\textbf{Goal:} Delete \texttt{VULKAN\_DISABLE\_WATERGEOM} environment variable check.

\textbf{Steps:}
\begin{enumerate}
    \item Remove \texttt{envDisableWaterGeom\_} member from WaterRenderer.hpp.
    \item Remove \texttt{std::getenv("VULKAN\_DISABLE\_WATERGEOM")} check in WaterRenderer::init (WaterRenderer.cpp:52).
    \item Remove all \texttt{if (envDisableWaterGeom\_)} guards in WaterRenderer::render, renderBrushLiquid.
    \item Delete \texttt{VULKAN\_DISABLE\_WATERGEOM=1} from any CI/scripts documentation.
    \item Verify: water pass always executes; no dead code paths remain.
\end{enumerate}

\textbf{Files:} \texttt{vulkan/renderer/WaterRenderer.hpp/cpp}
\end{tcolorbox}

\begin{tcolorbox}[promptbox={Prompt: Remove StagingRingBuffer Fallback from IndirectRenderer}]
\textbf{Goal:} Make UploadManager the only upload path; delete StagingRingBuffer + dedicated buffer fallback.

\textbf{Steps:}
\begin{enumerate}
    \item In \texttt{IndirectRenderer::uploadMeshes}, remove lines 622–658 (StagingRingBuffer allocation + fallback).
    \item Remove \texttt{pendingTransfer} struct and all related logic (publishPendingTransfer, pollPendingTransfers).
    \item Require \texttt{uploadMgr\_} to be non-null in \texttt{uploadMeshes} (assert or return false).
    \item Increase UploadManager slot count to 32, slot size to 4 MiB (covers largest chunk).
    \item Remove \texttt{StagingRingBuffer} dependency from IndirectRenderer; keep only in VulkanApp for texture/UBO uploads.
    \item Verify: \texttt{make run-debug} with GPU-assisted validation — no buffer barrier errors.
\end{enumerate}

\textbf{Files:} \texttt{vulkan/renderer/IndirectRenderer.hpp/cpp}, \texttt{vulkan/streaming/UploadManager.hpp/cpp}
\end{tcolorbox}

\begin{tcolorbox}[promptbox={Prompt: Remove Legacy Append Mode (Rebuild Path)}]
\textbf{Goal:} Delete legacy \texttt{addMesh}/\texttt{updateMesh}/\texttt{removeMesh}/\texttt{rebuild} API; only slotted mode remains.

\textbf{Steps:}
\begin{enumerate}
    \item In IndirectRenderer.hpp: remove legacy API declarations (lines 113–122).
    \item In IndirectRenderer.cpp: delete \texttt{addMesh}, \texttt{updateMesh}, \texttt{removeMesh}, \texttt{rebuild} implementations.
    \item Remove \texttt{mergedVertices}, \texttt{mergedIndices}, \texttt{indirectCommands} vectors (legacy CPU-side buffers).
    \item Remove \texttt{vertexBuffer}, \texttt{indexBuffer} single-buffer mirrors (slotted mode uses \texttt{vertexSlots}/\texttt{indexSlots}).
    \item Remove \texttt{metaBuffersWrittenCount}, \texttt{dirty}, \texttt{nextId} legacy tracking.
    \item Keep only slotted-mode members: \texttt{slotAlloc}, \texttt{spaceAlloc}, \texttt{slottedMode}, \texttt{meshes} map.
    \item Update SceneRenderer::initSlottedMode to be the only initialization path.
    \item Verify: no calls to legacy API remain in codebase (\texttt{grep -r "addMesh\\textbar rebuild" --include="*.cpp"}).
\end{enumerate}

\textbf{Files:} \texttt{vulkan/renderer/IndirectRenderer.hpp/cpp}, \texttt{vulkan/renderer/SceneRenderer.cpp}
\end{tcolorbox}

\begin{tcolorbox}[promptbox={Prompt: Remove Indirect Draw Count Fallback}]
\textbf{Goal:} Require \texttt{vkCmdDrawIndexedIndirectCount} (core since Vulkan 1.2); remove KHR function pointer loading.

\textbf{Steps:}
\begin{enumerate}
    \item In VulkanApp.cpp device creation: require \texttt{VK\_KHR\_draw\_indirect\_count} or Vulkan 1.2+.
    \item Remove \texttt{cmdDrawIndexedIndirectCount} function pointer from VegetationRenderer, DebugSDFRenderer.
    \item Replace all \texttt{cmdDrawIndexedIndirectCount} calls with direct \texttt{vkCmdDrawIndexedIndirectCount}.
    \item Remove error-log skip paths (DebugSDFRenderer.cpp:368–375).
    \item Update CMake/minimum Vulkan version to 1.2 if not already.
    \item Verify: builds on Vulkan 1.2+ devices; fails cleanly on 1.0/1.1 with clear message.
\end{enumerate}

\textbf{Files:} \texttt{vulkan/VulkanApp.cpp}, \texttt{vulkan/renderer/VegetationRenderer.cpp}, \texttt{vulkan/renderer/DebugSDFRenderer.cpp}
\end{tcolorbox}

\begin{tcolorbox}[promptbox={Prompt: Remove Descriptor Set Deferred Destruction Fallback}]
\textbf{Goal:} With Descriptor Buffers (VK\_EXT\_descriptor\_buffer), deferred set destruction becomes unnecessary.

\textbf{Steps:}
\begin{enumerate}
    \item After implementing Descriptor Buffer prompt (Section 7.2), remove all \texttt{deferDestroyUntilAllPending} calls for descriptor sets.
    \item Remove \texttt{vegDescriptorVersion}, \texttt{vegDescriptorSet} recreation logic in VegetationRenderer::onTextureArraysReallocated.
    \item Remove \texttt{sceneTextureWriteCache} in WaterRenderer (PostProcessRenderer already has similar).
    \item Descriptor buffer updates are simple memory writes — no set allocation/free needed.
    \item Verify: zero \texttt{vkFreeDescriptorSets} calls in render loop.
\end{enumerate}

\textbf{Files:} \texttt{vulkan/renderer/VegetationRenderer.cpp}, \texttt{vulkan/renderer/WaterRenderer.cpp}, \texttt{vulkan/renderer/PostProcessRenderer.cpp}
\end{tcolorbox}

\begin{tcolorbox}[promptbox={Prompt: Remove Shadow Pass Serial Fallback}]
\textbf{Goal:} Require parallel cascade rendering; remove serial \texttt{render()} path.

\textbf{Steps:}
\begin{enumerate}
    \item In ShadowRenderer::renderParallel, remove fallback block (lines 776–787).
    \item Require \texttt{cascadeSetsBuilt\_} to be true (assert).
    \item Remove \texttt{ShadowRenderer::render()} serial implementation entirely.
    \item Ensure \texttt{ensureShadowParallelResources} always succeeds (pre-create all resources at init).
    \item If device lacks multiple graphics queues, use single queue with semaphore chaining (still parallel CBs).
    \item Verify: \texttt{renderParallel} is only entry point; no serial shadow path remains.
\end{enumerate}

\textbf{Files:} \texttt{vulkan/renderer/ShadowRenderer.cpp}
\end{tcolorbox}

\begin{tcolorbox}[promptbox={Prompt: Remove Post-Process Sky Fallback}]
\textbf{Goal:} Sky offscreen targets always available; remove \texttt{skyView == VK\_NULL\_HANDLE} check.

\textbf{Steps:}
\begin{enumerate}
    \item In SceneRenderer::init, ensure SkyRenderer::createOffscreenTargets called before PostProcessRenderer::init.
    \item In PostProcessRenderer::render, remove lines 258–264 (skyView fallback to sceneColorView).
    \item Add assert \texttt{skyView != VK\_NULL\_HANDLE} at render entry.
    \item Verify: sky equirectangular targets always created (SkyRenderer.cpp:227+).
\end{enumerate}

\textbf{Files:} \texttt{vulkan/renderer/PostProcessRenderer.cpp}, \texttt{vulkan/renderer/SceneRenderer.cpp}
\end{tcolorbox}

\begin{tcolorbox}[promptbox={Prompt: Remove Cubemap Dummy Resources}]
\textbf{Goal:} Real cubemap always ready before first use; remove dummy 1×1 cubemap creation.

\textbf{Steps:}
\begin{enumerate}
    \item In Solid360Renderer::createSolid360Targets, ensure cubemap created before any consumer (WaterRenderer, PostProcessRenderer).
    \item Remove \texttt{cube360DummyColorImage}, \texttt{cube360DummyCubeView}, \texttt{cubemapDummyCubeView}, \texttt{cubemapDummyDepthView} members.
    \item Remove dummy creation in WaterRenderer::ensureCubemapResources (WaterRenderer.cpp:931–952).
    \item Remove dummy creation in Solid360Renderer::createSolid360Targets (lines 173–202).
    \item In WaterRenderer::updateSceneTexturesBinding, remove dummy fallback logic (lines 708–710).
    \item Verify: no dummy image views bound in any descriptor set at draw time.
\end{enumerate}

\textbf{Files:} \texttt{vulkan/renderer/WaterRenderer.cpp}, \texttt{vulkan/renderer/Solid360Renderer.cpp}
\end{tcolorbox}

\section{Conclusion}

The engine demonstrates a mature, modern Vulkan architecture with GPU-driven rendering as the dominant paradigm. The merged compute cull dispatch (solid + vegetation + SDF + bbox) is a significant optimization reducing dispatch overhead. Slotted-mode indirect rendering eliminates rebuild stalls. Primary bottlenecks are CPU-side descriptor management and staging buffer serialization — both addressable with VK\_EXT\_descriptor\_buffer and full UploadManager adoption. The detailed prompts above provide actionable starting points for each optimization and fallback removal.

\end{document}